\chapter{SPARQL}
\label{chap:sparql}
Im letzten Kapitel wurde aufgezeigt, wie Inferenz in einer Ontologie mittels der Regelsprache SWRL vollumfänglich genutzt werden kann.

Wie können Informationen aus der Ontologie und den darauf basierenden Inferenzen abgefragt werden?\\
Mit SPARQL --- einer Abfragesrprache --- ist es möglich.

Im folgenden Kapitel, welches, sofern nicht anders vermerkt, auf~\cite{w3sparql_querylang} basiert, stellen wir diese Abfragesprache vor.

Bei SPARQL handelt es sich sich um eine Abfragesprache für RDF.\@ Sie erlaubt es, Abfragen in mehreren Datenquellen vorzunehmen. Dabei werden Anfragen über Graphen vollzogen, auch entlang derer Konjunktionen und Disjunktionen. SPARQL unterstützt weiter Aggregation, Unterabfragen, Negation sowie die Nutzung von Ausdrücken als Werte.

Resultate sind entweder eine Menge von Ergebnissen oder RDF-Graphen.

(vgl.~\cite[Abstract]{w3sparql_querylang})

\noindent\rule[1ex]{\textwidth}{1pt}
\begin{wrapfigure}[7]{l}{0.1\textwidth}
    \vspace{-12pt}
    \includegraphics[width=0.1\textwidth]{bilder/owl.png}
\end{wrapfigure}
Kennt man die Abfragesprache SQL, so erscheint SPARQL auf den ersten Blick recht ähnlich --- der Name lässt dies bereits vermuten. Dies ist nur teilweise der Fall. Klauseln wie \textit{SELECT} und \textit{WHERE} sind ähnlich. In SPARQL gibt es jedoch keine \textit{FROM}-Klausel.\\
Der Hauptunterschied zwischen SPARQL und SQL ist jedoch die Verwendung von Variablen.  Diese werden in SPARQL in der Regel  bei jeder Abfrage verwendet. Abfragen ohne Variablen kommen in der Praxis (gemäss unserer Erfahrung) eher selten vor.\\
SQL enthält in der Regel sehr wenig Variablen.

\noindent\rule[1ex]{\textwidth}{1pt}

Eine SPARQL-Abfrage besteht aus folgenden Komponenten:
\begin{itemize}
    \item Einem oder mehreren Namespaces
    \item Variablen
    \item Einem (Teil-) Graphen
    \item Einer Menge von Gruppierungen und Aggregationen
    \item Einer Menge von Modifikatoren
    \item Abfragearten
    \item Ausdrücken und Wertevergleichen
\end{itemize}
(vgl.~\cite[18.1.10 SPARQL Query]{w3sparql_querylang})

\noindent\rule[1ex]{\textwidth}{1pt}
\begin{wrapfigure}[2]{l}{0.1\textwidth}
    \vspace{-19pt}
    \includegraphics[width=0.1\textwidth]{bilder/owl.png}
\end{wrapfigure}
Für die Praxis besteht eine SPARQL-Abfrage in der Regel aus \textit{Namespaces}, einer \textit{SELECT}-Klausel, einer \textit{WHERE}-Klausel sowie \textit{Modifikatoren}.\\

\noindent\rule[1ex]{\textwidth}{1pt}

\section{Beispiel einer SPARQL Abfrage}
\label{sec:sparql_beispiel}

Gegeben sei die folgende Datenbasis:
\lstset{caption={Einfache Datenbasis direkt in SPARQL\protect\footnotemark},captionpos=b}
\begin{lstlisting}
   <http://example.org/book/book1> <http://purl.org/dc/elements/1.1/title> "SPARQL Tutorial".
\end{lstlisting}
\footnotetext{\cite[2.1 Writing a Simple Query]{w3sparql_querylang}}

Die Datenbasis beinhaltet das Objekt \textit{book1} mit dem Attribut \textit{title}, welches den Wert \textit{``SPARQL Tutorial''} enthält.

Möchte ich nun den Titel des Buches abfragen, so kann ich folgende Abfrage verwenden:
\lstset{caption={Beispiel einer einfachen SPARQL-Abfrage\protect\footnotemark},captionpos=b}
\begin{lstlisting}
    SELECT
        ?title
    WHERE {
        <http://example.org/book/book1> <http://purl.org/dc/elements/1.1/title> ?title .
    }
\end{lstlisting}
\footnotetext{\cite[2.1 Writing a Simple Query]{w3sparql_querylang}}

Sie ergibt folgendes Resultat:
\noindent\hspace*{15mm}
\begin{table}[h]
    \centering
    \begin{tabular}{|l|}
        \hline
        \multicolumn{1}{|c|}{\textbf{title}} \\ \hline
        ``SPARQL Tutorial''                    \\ \hline
    \end{tabular}
    \caption{Resultat einer einfachen SPARQL-Abfrage\protect\footnotemark}
\end{table}
\footnotetext{\cite[2.1 Writing a Simple Query]{w3sparql_querylang}}

\section{Namespaces}
\label{sec:sparql_namespaces}

SPARQL ist eine Abfragesprach für RDF, das seinerseits auf XML basiert. Dadurch sind XML-Namespaces auch in SPARQL verfügbar. Bei Namespaces handelt es sich um Referenzen auf andere XML-basierte Dokumente, welche modular genutzt werden können.\\
Ein Namespace besteht aus einem Kürzel (welches auch leer sein kann), sowie aus der kompletten Adresse der Ressource (vgl.~\cite[2.1 Introduction]{w3rdf_syntax}).

\lstset{language=XML,
  caption={Beispiel von Namespaces in RDF\protect\footnotemark},
  frame=L,
  basicstyle=\small\normalfont\sffamily,  % the size of the fonts that are used for the code
  stepnumber=1,                           % the step between two line-numbers.
                                          % If it is 1 each line will be numbered
  numbersep=10pt,                         % how far the line-numbers are from the code
  tabsize=2,                              % tab size in blank spaces
  extendedchars=true,                     %
  breaklines=true,                        % sets automatic line breaking
  captionpos=b,                           % sets the caption-position to bottom
  mathescape=true,
  showspaces=false,                       % Leerzeichen anzeigen ?
  showtabs=false,                         % Tabs anzeigen ?
  xleftmargin=17pt,
  framexleftmargin=17pt,
  framexrightmargin=17pt,
  framexbottommargin=5pt,
  framextopmargin=5pt,
  showstringspaces=false                  % Leerzeichen in Strings anzeigen ?
  belowcaptionskip=5em,
  belowskip=3em,
  aboveskip=3em
 }
\begin{lstlisting}
$\textless$rdf:RDF xmlns:rdf="http://www.w3.org/1999/02/22-rdf-syntax-ns#"  % chktex 1
         xmlns:dc="http://purl.org/dc/elements/1.1/"
         xmlns:ex="http://example.org/stuff/1.0/"$\textgreater$
\end{lstlisting}
\footnotetext{\cite[2.6 Completing the Document: Document Element and XML Declaration]{w3rdf_syntax}}

\section{Variablen}
\label{sec:sparql_variablen}
Bei Abfragen unterstützt SPARQL die Verwendung von Variablen. Diese können dabei an einer beliebigen Stelle, wie bei (Teil-) Graphen, Gruppierungen, Aggregationen und Modifikatoren eingesetzt werden.\\
Die Variablen dienen als Platzhalter für abzufragende Objekte, Subjekte und Prädikate. Variablen beginnen entweder mit einem Dollarzeichen \textit{\$} oder einem Fragezeichen \textit{?}, welche selbst nicht Teil der Variable sind.

\lstset{caption={Beispiel einer SPARQL Abfrage mit den Variablen \textit{?object}, \textit{?predicate} und \textit{?subject}.},captionpos=b}
\begin{lstlisting}
    SELECT DISTINCT
        *
    WHERE {
      ?object ?predicate ?subject
    }
    LIMIT
        10
\end{lstlisting}

Wird obige Anfrage auf einen SPARQL-Endpunkt, wie z.B. \textit{http://dbpedia.org/}, angewendet, so ergibt dies folgendes Resultat:
\noindent\hspace*{15mm}
\begin{table}[h]
    \centering
    \begin{tabular}{l|l|l|}
        \hline
        \multicolumn{1}{|l|}{\textbf{object}}                                & \textbf{predicate}                               & \textbf{subject}                                  \\ \hline
        \multicolumn{1}{|l|}{http://dbpedia.org/ontology/acceleration}      & ... rdf-syntax-ns\#type & ... owl\#FunctionalProperty \\ \hline
        \multicolumn{1}{|l|}{http://dbpedia.org/ontology/torqueOutput}      & ... rdf-syntax-ns\#type & ... owl\#FunctionalProperty \\ \hline
        \multicolumn{1}{|l|}{http://dbpedia.org/ontology/birthDate}         & ... rdf-syntax-ns\#type & ... owl\#FunctionalProperty \\ \hline
        \multicolumn{1}{|l|}{http://dbpedia.org/ontology/birthYear}         & ... rdf-syntax-ns\#type & ... owl\#FunctionalProperty \\ \hline
        \multicolumn{1}{|l|}{http://dbpedia.org/ontology/deathDate}         & ... rdf-syntax-ns\#type & ... owl\#FunctionalProperty \\ \hline
        \multicolumn{1}{|l|}{http://dbpedia.org/ontology/deathYear}         & ... rdf-syntax-ns\#type & ... owl\#FunctionalProperty \\ \hline
        \multicolumn{1}{|l|}{http://dbpedia.org/ontology/diameter}          & ... rdf-syntax-ns\#type & ... owl\#FunctionalProperty \\ \hline
        \multicolumn{1}{|l|}{http://dbpedia.org/ontology/displacement}      & ... rdf-syntax-ns\#type & ... owl\#FunctionalProperty \\ \hline
        \multicolumn{1}{|l|}{http://dbpedia.org/ontology/height}            & ... rdf-syntax-ns\#type & ... owl\#FunctionalProperty \\ \hline
        \multicolumn{1}{|l|}{http://dbpedia.org/ontology/latestReleaseDate} & ... rdf-syntax-ns\#type & ... owl\#FunctionalProperty \\ \hline
    \end{tabular}
    \caption{Resultat einer Abfrage des DBPedia-Endpunktes limitiert auf 10 Ergebnisse.}
\end{table}

\section{(Teil-) Graphen}
\label{sec:sparql_graph}

SPARQL basiert auf dem Matching eines (Teil-) Graphen in einem gegebenen Graphen. Dabei ist ein (Teil-) Graph eine Menge von Tripeln.

Man unterscheidet zwischen den folgenden Arten von Teilgraphen:
\begin{itemize}
\item \textit{Grundlegende Teilgraphen}

Hierbei muss eine Menge von Tripeln derjenigen im gegebenen Graphen entsprechen

\item \textit{Gruppierende Teilgraphen}

Hierbei müssen \textit{alle} Elemente einer Menge von Tripeln von grundlegenden Teilgraphen wiederum derjenigen im gegebenen Graphen entsprechen

\item \textit{Optionale Teilgraphen}

Zusätzliche Teilgraphen können die Lösungsmenge erweitern

\item \textit{Alternative Teilgraphen}

Versuch, ob zwei oder mehrere mögliche Teilgraphen denjenigen im gegebenen Graphen entsprechen

\item \textit{Teilgraphen in namensbasierten Graphen}

Versuch, ob Teilgraphen denjenigen im gegebenen, namensbasierten Graphen entsprechen
\end{itemize}

Ein grundlegender Teilgraph ist somit eine Menge von Tripeln, wobei jeder Tripel gemäss RDF-Spezifikation immer aus Subjekt, Prädikat und Objekt bestehen muss.~\cite[3.1 Triples]{w3rdf}

\lstset{caption={Beispiel eines grundlegenden Teilgraphen in SPARQL}}
\begin{lstlisting}
    ?subject ?predicate ?object.
\end{lstlisting}

\noindent\rule[1ex]{\textwidth}{1pt}
\begin{wrapfigure}[3]{l}{0.1\textwidth}
    \vspace{-16pt}
    \includegraphics[width=0.1\textwidth]{bilder/owl.png}
\end{wrapfigure}
Ein Teilgraph wird in SPARQL immer zur Gruppierung einer Abfrage verwendet. Dies entspricht in der SQL-Sprache der \textit{WHERE}-Klausel.\\
Ein Teilgraph wird immer mit dem Punkt-Operator beendet.

\noindent\rule[1ex]{\textwidth}{1pt}

Ein gruppierender Teilgraph wird in der SPARQL-Abfrag immer mit geschweiften Klammern \textit{\{\}} umschlossen. Es ist möglich die Resultate einer Gruppierung zu filtern. Dazu wird \textit{FILTER} als Schlüsselwort verwendet. Dieses bezieht sich immer auf die Gruppierung, in welcher es steht.

\lstset{caption={Beispiel eines gruppierenden Teilgraphen mit dem \textit{FILTER}-Schlüsselwort\protect\footnotemark}}
\begin{lstlisting}
    {
        ?x foaf:name ?name.
        ?x foaf:mbox ?mbox.
        FILTER regex(?name, "Smith")
    }
\end{lstlisting}
\footnotetext{\cite[5.2.2 Scope of Filters]{w3sparql_querylang}}

Im obigem Beispiel werden alle Objekte und die Werte ihrer Attribute (\textit{name} und \textit{mbox}, welche im Namespace \textit{foaf} bezeichnet werden) abgefragt.

Dabei kann eine komplette leere Menge --- wenn keine Objekte vorhanden sind --- oder aber eine leere Teilmenge --- wenn ein Objekt z.B. nicht über die beiden Attribute \textit{name} und \textit{mbox} verfügt --- zurückgegeben werden.

Es bleiben nur Objekte zurück, welche über das Attribut \textit{foaf:name} verfügen und dessen Inhalt \textit{``Smith''} ist.

\section{Gruppierungen und Aggregationen}
\label{sec:sparql_gruppierungenaggregationen}
Aggregationen bezeichnen eine bestimmte Art der Verbindung zwischen Objekten.~\cite{wiki:aggregation} Aggregationen wenden Ausdrücke bzw. Funktionen auf eine Lösungsmenge an. Dabei enthält eine Lösung normalerweise eine einzelne Gruppe mit allen Lösungen.

Gruppierungen werden mittels \textit{GROUP BY} angegeben.

SPARQL unterstützt aktuell die Aggregationen \textit{COUNT}, \textit{SUM}, \textit{MIN}, \textit{MAX}, \textit{AVG}, \textit{GROUP\_CONCAT} und \textit{SAMPLE}.

Die Aggregation wird dann verwendet, wenn ein Resultat über eine Gruppe von Lösungen berechnet werden soll (Beispiel: Maximaler Wert einer bestimmten Variablen).

(vgl.~\cite[11 Aggregates]{w3sparql_querylang})

%%%%% \lstset{caption={Beispiel einer Abfrage mit Gruppierung und Aggregation\protect\footnotemark}}
%%%%% \lstset{language=SQL}
%%%%% \begin{lstlisting}
%%%%%     PREFIX books: $\textless$http://books.example/$\textgreater$
%%%%%     SELECT (
%%%%%         SUM(?lprice) AS ?totalPrice
%%%%%     )
%%%%%     WHERE {
%%%%%       ?org books:affiliates ?auth .
%%%%%       ?auth books:writesBook ?book .
%%%%%       ?book books:price ?lprice .
%%%%%     }
%%%%%     GROUP BY ?org
%%%%%     HAVING (SUM(?lprice) $\textgreater$ 10)
%%%%% \end{lstlisting}
%%%%% \footnotetext{\cite[11.1 Aggregate Example]{w3sparql_querylang}}
%%%%%     
%%%%%     \section{Modifikatoren}
%%%%%     \label{sec:sparql_modifikatoren}
%%%%%     
%%%%%     Abfragen mittels SPARQL generieren eine ungeordnete Menge von Lösungen, wobei jede dieser Lösungen eine Funktion von Variablen zu RDF-Termen darstellt. Solche Lösungen werden als Sequenz von Lösungen behandelt, initial allerdings ohne spezifische Ordnung.~\cite[15 Solution Sequences and Modifiers]{w3sparql_querylang}
%%%%%     
%%%%%     Um eine Sequenz von Lösungen zu ordnen, können Modifikatoren verwendet werden. Man unterscheidet dabei zwischen folgenden Modifikatoren:
%%%%%     
%%%%%     \begin{itemize}
%%%%%     \item \textit{Order}
%%%%%     
%%%%%     Sortiert die Lösungen nach einer gewissen Variable, entweder ab- oder aufsteigend
%%%%%     
%%%%%      \item \textit{Projection}
%%%%%      
%%%%%      Erlaubt die Auswahl von spezifischen Variablen mittels der \textit{Select}-Klausel
%%%%%      
%%%%%      \item \textit{Distinct}
%%%%%      
%%%%%      Stellt sicher, dass die Lösungen in der Lösungsmenge einmalig sind
%%%%%      
%%%%%      \item \textit{Reduced}
%%%%%      
%%%%%      Verhindert die Elimination von Duplikaten in der Lösungsmenge
%%%%%      
%%%%%      \item \textit{Offset}
%%%%%      
%%%%%      Definiert ab welcher Lösung die gefundenen Lösungen ausgegeben werden sollen
%%%%%      
%%%%%      \item \textit{Limit}
%%%%%      
%%%%%      Definiert die maximale Anzahl der gefunden Lösungen, welche ausgegeben werden sollen
%%%%%      
%%%%%      \end{itemize}
%%%%%      
%%%%%      Die Modifikatoren werden in der Reihenfolge der obigen Liste angewendet.~\cite[vgl. 15 Solution Sequences and Modifiers]{w3sparql_querylang}
%%%%%      
%%%%%      
%%%%%      \section{Abfragearten}
%%%%%      \label{sec:sparql_abfragearten}
%%%%%      
%%%%%      SPARQL unterscheidet zwischen vier Abfragearten:
%%%%%      \begin{itemize}
%%%%%          \item \textit{SELECT}
%%%%%      
%%%%%              Gibt alle bzw.\ eine Teilmenge der Variablen eines (Teil-) Graphen zurück
%%%%%          \item \textit{CONSTRUCT}
%%%%%      
%%%%%              Gibt einen RDF Graphen durch Substitution der Variablen in einer Menge von gegebenen Tripeln zurück
%%%%%          \item \textit{ASK}
%%%%%      
%%%%%              Gibt einen Wahrheitswert (Boolean) zurück, welcher angibt, ob die Menge von angefragten Tripeln denen in dem gegebenen Graphen entsprechen oder nicht
%%%%%          \item \textit{DESCRIBE}
%%%%%      
%%%%%              Gibt einen RDF Graphen zurück, welcher die gefundenen Ressourcen beschreibt
%%%%%      \end{itemize}
%%%%%      
%%%%%      \subsection{SELECT}
%%%%%      \label{subsec:sparql_abfragearten_select}
%%%%%      Die \textit{SELECT}-Abfrageart gibt Variablen und deren Belegung direkt, in Form von Projektionen, zurück. Dabei werden neue Variablen-Belegungen generiert. Die spezifischen, gewünschten Variablen werden in Form einer Liste, per Leerschlag getrennt, angegeben.
%%%%%      
%%%%%      \lstset{caption={Beispiel der \textit{SELECT}-Abfrageart in SPARQL},captionpos=b,language=SQL}
%%%%%      \begin{lstlisting}
%%%%%      SELECT ?a ?b ?c
%%%%%      WHERE {
%%%%%      ...
%%%%%      \end{lstlisting}
%%%%%      
%%%%%      Die $ SELECT * $ Syntax ist dabei eine Kurzschreibweise und bedeutet eigentlich die Verwendung bzw. Auflistung aller in der Abfrage verwendeten Variablen. Ausgenommen sind Variablen innerhalb von \textit{FILTER}-Anweisungen sowie Variablen, welche rechts des Minus-Operators stehen. Die Syntax darf nur verwendet werden, wenn die Abfrage keine \textit{GROUP BY} Aggregation verwendet.
%%%%%      \lstset{caption={Beispiel der \textit{SELECT *} Kurzschreibweise in SPARQL},captionpos=b,language=SQL}
%%%%%      \begin{lstlisting}
%%%%%      SELECT *
%%%%%      WHERE {
%%%%%      ...
%%%%%      \end{lstlisting}
%%%%%      
%%%%%      Die \textit{SELECT}-Abfrageart bietet die Möglichkeit neue Variable durch Ausdrücke einzuführen. Diese werden mittels $(expr AS V)$ eingeführt.
%%%%%      
%%%%%      \lstset{caption={Beispiel eines Ausdrucks der \textit{SELECT}-Abfrageart},captionpos=b,language=SQL}
%%%%%      \begin{lstlisting}
%%%%%      SELECT ?title (?p * (1 - ?discount) AS ?price)
%%%%%      WHERE {
%%%%%      ...
%%%%%      \end{lstlisting}
%%%%%      
%%%%%      
%%%%%      \subsection{CONSTRUCT}
%%%%%      \label{subsec:sparql_abfragearten_construct}
%%%%%      Die \textit{CONSTRUCT}-Abfrageart gibt einen RDF Graphen durch Substitution der Variablen in einer Menge von gegebenen Tripeln zurück. Dabei wird jede Lösung der Abfrage in die Lösungssequenz, unter Substitution der Variablen in der gegebenen Vorlage des Graphen, überführt. Die Lösungsmenge wird schliesslich per Vereinigung auf einen einzigen RDF Graphen reduziert.
%%%%%      
%%%%%      \lstset{caption={Beispiel einer SPARQL Abfrage mit Nutzung der \textit{CONSTRUCT}-Abfrageart\protect\footnotemark},captionpos=b}
%%%%%      \begin{lstlisting}
%%%%%      PREFIX foaf:    $\textless$http://xmlns.com/foaf/0.1/$\textgreater$
%%%%%      PREFIX vcard:   $\textless$http://www.w3.org/2001/vcard-rdf/3.0#$\textgreater$
%%%%%      CONSTRUCT   { $\textless$http://example.org/person#Alice$\textgreater$ vcard:FN ?name }
%%%%%      WHERE       { ?x foaf:name ?name }
%%%%%      \end{lstlisting}
%%%%%      \footnotetext{\cite[16.2 CONSTRUCT]{w3sparql_querylang}}
%%%%%      
%%%%%      Die obige Abfrage erzeugt folgende Ausgabe:
%%%%%      \lstset{caption={Beispiel der Ausgabe einer SPARQL Abfrage mit Nutzung der \textit{CONSTRUCT}-Abfrageart\protect\footnotemark},captionpos=b}
%%%%%      \begin{lstlisting}
%%%%%      @prefix vcard: $\textless$http://www.w3.org/2001/vcard-rdf/3.0#$\textgreater$ .
%%%%%      
%%%%%      $\textless$http://example.org/person#Alice$\textgreater$ vcard:FN "Alice" .
%%%%%      \end{lstlisting}
%%%%%      \footnotetext{\cite[16.2 CONSTRUCT]{w3sparql_querylang}}
%%%%%      
%%%%%      Wird bei diesem Vorgang ein ungültiges Tripel erzeugt, so wird dieses Tripel im Ausgabegraphen nicht angezeigt. Enthält die Vorlage eine Graphen Tripel ohne Variablen, erscheinen diese dennoch im Ausgabegraphen.
%%%%%      
%%%%%      \lstset{caption={Beispiel einer SPARQL Abfrage mit Nutzung der \textit{CONSTRUCT}-Abfrageart unter Verwendung leerer Knoten\protect\footnotemark},captionpos=b}
%%%%%      \begin{lstlisting}
%%%%%      PREFIX foaf:    $\textless$http://xmlns.com/foaf/0.1/$\textgreater$
%%%%%      PREFIX vcard:   $\textless$http://www.w3.org/2001/vcard-rdf/3.0#$\textgreater$
%%%%%      
%%%%%      CONSTRUCT { ?x  vcard:N _:v .
%%%%%                  _:v vcard:givenName ?gname .
%%%%%                  _:v vcard:familyName ?fname }
%%%%%      WHERE
%%%%%       {
%%%%%          { ?x foaf:firstname ?gname } UNION  { ?x foaf:givenname   ?gname } .
%%%%%          { ?x foaf:surname   ?fname } UNION  { ?x foaf:family_name ?fname } .
%%%%%       }
%%%%%      \end{lstlisting}
%%%%%      \footnotetext{\cite[16.2.1 Templates with Blank Nodes]{w3sparql_querylang}}
%%%%%      
%%%%%      Die obige Abfrage erzeugt folgende Ausgabe:
%%%%%      \lstset{caption={Beispiel der Ausgabe einer SPARQL Abfrage mit Nutzung der \textit{CONSTRUCT}-Abfrageart unter Verwendung leerer Knoten\protect\footnotemark},captionpos=b}
%%%%%      \begin{lstlisting}
%%%%%      @prefix vcard: $\textless$http://www.w3.org/2001/vcard-rdf/3.0#$\textgreater$ .
%%%%%      
%%%%%      _:v1 vcard:N         _:x .
%%%%%      _:x vcard:givenName  "Alice" .
%%%%%      _:x vcard:familyName "Hacker" .
%%%%%      
%%%%%      _:v2 vcard:N         _:z .
%%%%%      _:z vcard:givenName  "Bob" .
%%%%%      _:z vcard:familyName "Hacker" .
%%%%%      \end{lstlisting}
%%%%%      \footnotetext{\cite[16.2.1 Templates with Blank Nodes]{w3sparql_querylang}}
%%%%%      
%%%%%      \subsection{ASK}
%%%%%      \label{subsec:sparql_abfragearten_ask}
%%%%%      Die \textit{ASK}-Abfrageart ermöglicht die Ausgabe, ob eine Lösung für die gegebene Abfrage existiert oder nicht.
%%%%%      
%%%%%      Gegeben sei die folgende Datenbasis:
%%%%%      \lstset{caption={Einfache Datenbasis direkt in SPARQL\protect\footnotemark},captionpos=b}
%%%%%      \lstset{language=XML}
%%%%%      \begin{lstlisting}
%%%%%      @prefix foaf:       $\textless$http://xmlns.com/foaf/0.1/$\textgreater$ .
%%%%%      
%%%%%      _:a  foaf:name       "Alice" .
%%%%%      _:a  foaf:homepage   $\textless$http://work.example.org/alice/$\textgreater$ .
%%%%%      
%%%%%      _:b  foaf:name       "Bob" .
%%%%%      _:b  foaf:mbox       $\textless$mailto:bob@work.example$\textgreater$ .
%%%%%      \end{lstlisting}
%%%%%      \footnotetext{\cite[16.3 ASK]{w3sparql_querylang}}
%%%%%      
%%%%%      Wird nun folgende Abfrage getätigt:
%%%%%      \lstset{caption={Beispiel einer SPARQL-Abfrage mit Nutzung der \textit{ASK}-Abfrageart\protect\footnotemark},captionpos=b}
%%%%%      \lstset{language=XML}
%%%%%      \begin{lstlisting}
%%%%%      PREFIX foaf:    $\textless$http://xmlns.com/foaf/0.1/$\textgreater$
%%%%%      ASK  { ?x foaf:name  "Alice" }
%%%%%      \end{lstlisting}
%%%%%      \footnotetext{\cite[16.3 ASK]{w3sparql_querylang}}
%%%%%      So gibt das System $true.$ zurück, da die Datenbasis ein Individuum mit dem Attribut $foaf:name$ dessen Wert $'Alice'$ ist.
%%%%%      
%%%%%      \subsection{DESCRIBE}
%%%%%      \label{subsec:sparql_abfragearten_describe}
%%%%%      Die \textit{DESCRIBE}-Abfrageart gibt einen RDF Graphen mit Informationen in Form von RDF-Daten über Ressourcen zurück. Ressourcen können dabei \textit{IRI}s oder Variablen sein.
%%%%%      
%%%%%      \lstset{caption={Beispiel der \textit{DESCRIBE}-Abfrageart in SPARQL},captionpos=b,language=SQL}
%%%%%      \begin{lstlisting}
%%%%%      DESCRIBE ?a ?b ?c
%%%%%      WHERE {
%%%%%      ...
%%%%%      \end{lstlisting}
%%%%%      
%%%%%      Die $ DESCRIBE * $ Syntax ist dabei eine Kurzschreibweise und bedeutet eigentlich die Verwendung bzw. Auflistung aller in der Abfrage verwendeten Variablen. 
%%%%%      \lstset{caption={Beispiel der \textit{DESCRIBE *} Kurzschreibweise in SPARQL},captionpos=b,language=SQL}
%%%%%      \begin{lstlisting}
%%%%%      DESCRIBE *
%%%%%      WHERE {
%%%%%      ...
%%%%%      \end{lstlisting}
%%%%%      
%%%%%      \section{Ausdrücke und Wertevergleiche}
%%%%%      \label{sec:sparql_ausdruecke}
%%%%%      
%%%%%      Jede Abfrage in SPARQL kann mit einem so genannten Filter versehen werden. Dabei wird versucht ein gewisses Muster, gegeben durch Beschränkungen, auf (Teil-) Graphen anzuwenden um so die (Teil-) Graphen einzuschränken oder eben zu filtern.
%%%%%      
%%%%%      Jedes RDF-Literal kann durch einen bestimmten Datentyp definiert sein, so zum Beispiel Wahrheitswerte (Boolean) oder Datum und Uhrzeit (DateTime). Der Filter in SPARQL erlaubt dabei den Wertevergleich auf typisierte RDF-Literale. Die dabei verwendeten Operanden und Datentypen finden sich unter~\footnote{http://www.w3.org/TR/xpath-functions/} und~\footnote{http://www.w3.org/TR/sparql11-query/\#operandDataTypes}.
%%%%%      
%%%%%      Eine SPARQL-Abfrage unter Verwendung eines Filters kann zum Beispiel wie folgt aussehen:
%%%%%      \lstset{caption={Beispiel einer einfachen SPARQL Abfrage unter Verwendung eines Filters},captionpos=b}
%%%%%      \begin{lstlisting}
%%%%%          ...
%%%%%          SELECT ?book
%%%%%          WHERE
%%%%%          {
%%%%%            ?book :isInLibrary :BerneUniversityLibrary
%%%%%            ?book :hasAmount ?amount
%%%%%            FILTER (?amount $\textgreater$ xsd:integer(10))
%%%%%          }
%%%%%      \end{lstlisting}
%%%%%      Hier werden alle Bücher abgefragt, welche sich in der Bibliothek der Universität von Bern befinden und deren (allgemeines) Vorkommen grösser als 10 ist.
%%%%%      
%%%%%      Ist man in obigem Beispiel nicht sicher, welcher Datentyp die Variable $?amount$ entspricht, so kann man die Funktionalität zur Umwandlung (Casting) durch SPARQL nutzen:
%%%%%      \lstset{caption={Beispiel einer einfachen SPARQL Abfrage unter Verwendung eines Filters mit Typenumwandlung},captionpos=b}
%%%%%      \begin{lstlisting}
%%%%%          ...
%%%%%          SELECT ?book
%%%%%          WHERE
%%%%%          {
%%%%%            ?book :isInLibrary :BerneUniversityLibrary
%%%%%            ?book :hasAmount ?amount
%%%%%            FILTER (xsd:integer(?amount) $\textgreater$ xsd:integer(10))
%%%%%          }
%%%%%      \end{lstlisting}
%%%%%      
%%%%%      Die genaue Auswertung eines Filters lässt sich unter~\footnote{http://www.w3.org/TR/sparql11-query/\#evaluation} nachlesen. Die interne Verwendung und Auswertung von Operatoren lässt sich unter~\footnote{http://www.w3.org/TR/sparql11-query/\#OperatorMapping} nachlesen.
%%%%%      
%%%%%      \subsection{Funktionen}
%%%%%      \label{subsec:sparql_ausdruecke_funktionen}
%%%%%      
%%%%%      SPARQL bietet eine Vielzahl an Operatoren und Funktionen. Diese alle einzeln aufzuführen und zu beschreiben würde den Rahmen dieses Dokumentes sprengen, daher werden nur die von den Autoren als am häufigsten genutzten Operatoren und Funktionen aufgeführt. Eine genauere Beschreibung findet sich unter~\footnote{http://www.w3.org/TR/sparql11-query/\#SparqlOps}.
%%%%%      
%%%%%      \subsubsection{BOUND}
%%%%%      \label{subsec:sparql_ausdruecke_funktionen_bound}
%%%%%      Mittels \textit{BOUND} lässt sich prüfen, ob ein Objekt an einen Wert gebunden ist oder nicht. Man kann dies als eine Art Überprüfung auf gewisse Eigenschaften verstehen.
%%%%%      
%%%%%      Möchte man zum Beispiel wissen, welche Personen alle über eine Relation zu einem Datum (z.B. Erfassungsdatum) verfügen, so geschieht dies mittels:
%%%%%      
%%%%%      \lstset{caption={Beispiel einer simplen Ontologie zur Nutzung der \textit{BOUND}-Funktion in einer Abfrage\protect\footnotemark},captionpos=b}
%%%%%      \begin{lstlisting}
%%%%%      @prefix foaf:        $\textless$http://xmlns.com/foaf/0.1/$\textgreater$ .
%%%%%      @prefix dc:          $\textless$http://purl.org/dc/elements/1.1/$\textgreater$ .
%%%%%      @prefix xsd:          $\textless$http://www.w3.org/2001/XMLSchema#$\textgreater$ .
%%%%%      
%%%%%      _:a  foaf:givenName  "Alice".
%%%%%      
%%%%%      _:b  foaf:givenName  "Bob" .
%%%%%      _:b  dc:date         "2005-04-04T04:04:04Z"^^xsd:dateTime .
%%%%%      \end{lstlisting}
%%%%%      \footnotetext{\cite[17.4.1.1 BOUND]{w3sparql_querylang}}
%%%%%      
%%%%%      \lstset{caption={Beispiel einer einfachen SPARQL Abfrage zur Nutzung der \textit{BOUND}-Funktion\protect\footnotemark},captionpos=b}
%%%%%      \begin{lstlisting}
%%%%%      PREFIX foaf: $\textless$http://xmlns.com/foaf/0.1/$\textgreater$
%%%%%      PREFIX dc:   $\textless$http://purl.org/dc/elements/1.1/$\textgreater$
%%%%%      PREFIX xsd:   $\textless$http://www.w3.org/2001/XMLSchema#$\textgreater$
%%%%%      SELECT ?givenName
%%%%%       WHERE { ?x foaf:givenName  ?givenName .
%%%%%               OPTIONAL { ?x dc:date ?date } .
%%%%%               FILTER ( bound(?date) ) }
%%%%%      \end{lstlisting}
%%%%%      \footnotetext{\cite[17.4.1.1 BOUND]{w3sparql_querylang}}
%%%%%      
%%%%%      Die obige Abfrage ergibt folgendes Resultat:
%%%%%      \noindent\hspace*{15mm}
%%%%%      \begin{table}[h]
%%%%%          \centering
%%%%%          \begin{tabular}{|l|}
%%%%%              \hline
%%%%%              \multicolumn{1}{|c|}{\textbf{givenName}} \\ \hline
%%%%%              ``Bob''                    \\ \hline
%%%%%          \end{tabular}
%%%%%          \caption{Resultat einer SPARQL Abfrage mittels \textit{BOUND}-Filter\protect\footnotemark}
%%%%%      \end{table}
%%%%%      \footnotetext{\cite[17.4.1.1 BOUND]{w3sparql_querylang}}
%%%%%      
%%%%%      Es wird also nur die Person \textit{``Bob''} gefunden, da nur diese Person gemäss definierter Ontologie über eine Relation zu einem Datum verfügt.
%%%%%      
%%%%%      \subsubsection{IF}
%%%%%      \label{subsec:sparql_ausdruecke_funktionen_if}
%%%%%      Die \textit{IF}-Funktion benötigt drei Argumente als Parameter. Dabei fasst sie das erste Argument als Wahrheitswert zur Auswertung auf. Das zweite Argument wird zurückgegeben, wenn der Ausdruck des ersten Argumentes wahr ist. Das dritte Argument wird zurückgegeben, wenn der Ausdruck des ersten Argumentes falsch ist.
%%%%%      
%%%%%      \lstset{caption={Beispiel der \textit{IF}-Funktion in SPARQL},captionpos=b,language=XML}
%%%%%      \begin{lstlisting}
%%%%%      IF(?a = 2,?b,?c)
%%%%%      \end{lstlisting}
%%%%%      
%%%%%      Oder, informal, in prozeduraler Programmierung ausgedrückt:
%%%%%      \lstset{caption={Beispiel einer simplen Ontologie zur Nutzung der \textit{BOUND}-Funktion in einer Abfrage\protect\footnotemark},captionpos=b}
%%%%%      \lstset{caption={Beispiel der \textit{IF}-Funktion in SPARQL anhand prozeduraler Programmierung},
%%%%%          captionpos=b,
%%%%%          language=Java
%%%%%      }
%%%%%      \begin{lstlisting}
%%%%%      if (a == 2) {
%%%%%          return b;
%%%%%      } else {
%%%%%          return c;
%%%%%      }
%%%%%      \end{lstlisting}
%%%%%      
%%%%%      \subsubsection{IN}
%%%%%      \label{subsec:sparql_ausdruecke_funktionen_in}
%%%%%      Die \textit{IN}-Funktion teste, ob der Parameter, welcher direkt vor ihrem Aufruf folgt, in dem Listenausdruck vorkommt, welche direkt nach ihrem Aufruf folgt. Ist dies der Fall, so gibt die Funktion den Wahrheitswert wahr zurück, ansonsten falsch.
%%%%%      
%%%%%      \lstset{caption={Beispiel der \textit{IN}-Funktion in SPARQL},captionpos=b,language=SQL}
%%%%%      \begin{lstlisting}
%%%%%      2 IN (1, 2, 3) % Wahr
%%%%%      2 IN ()        % Falsch
%%%%%      \end{lstlisting}
%%%%%      \footnotetext{\cite[17.4.1.9 IN]{w3sparql_querylang}}
%%%%%      In den obigen Beispielen wäre der erste Ausdruck also wahr, da 2 in der gegebenen Liste von (1, 2, 3) vorkommt, der zweite Ausdruck falsch, da 2 nicht in der leeren Menge vorkommt.
%%%%%      
%%%%%      Analog zur \textit{IN}-Funktion existiert die \textit{NOT IN}-Funktion. Es handelt sich dabei um die Negation der \textit{IN}-Funktion.
%%%%%      
%%%%%      
%%%%%      \noindent\rule[1ex]{\textwidth}{1pt}
%%%%%      \begin{wrapfigure}[7]{l}{0.1\textwidth}
%%%%%          \vspace{-12pt}
%%%%%          \includegraphics[width=0.1\textwidth]{bilder/elephant.png}
%%%%%      \end{wrapfigure}
%%%%%      Wie bereits beschrieben, sind alle Komponenten vorhanden, um Ausflüge abbilden zu können, welche den unter~\hyperref[elephant_graph_data]{\nameref{chap:graph_data}} genannten Kriterien entsprechen. Mit diesem Kapitel ist nun auch noch die fehlende Abfragesprache um Abfragen in der Ontologie vornehmen zu können eingeführt worden.
%%%%%      
%%%%%      Möchte man nun, um auf das beschriebene Beispiel der Familie, welche einen Ausflug planen möchte, zurückzukommen, alle Ausflüge der Ontologie ausgeben, welche den Kriterien  \textit{familienfreundlich}, \textit{regional} und \textit{actionreich} entsprechen, so gelingt dies mit der nacholgenden Abfrage.
%%%%%      
%%%%%      \begin{lstlisting}[caption={Beispiel einer Abfrage um alle familiengerechten, actionreichen und regionalen Ausflüge der Ontologie auszugeben},captionpos=b,language=SQL]
%%%%%          SELECT
%%%%%              *
%%%%%          WHERE {
%%%%%              ?object :familiengerecht true;
%%%%%                  :regional true;
%%%%%                  :action true.
%%%%%          }
%%%%%      \end{lstlisting}
%%%%%      
%%%%%      \noindent\rule[1ex]{\textwidth}{1pt}
